%----------------------------------------------------------------------%
%  Plan de Trading (v.3 – Revisión corregida)                         %
%----------------------------------------------------------------------%

\chapter{Plan de Trading}

\begin{itemize}
\item El \textbf{trading es un trabajo} y, como cualquier otro negocio, \textbf{requiere un plan}.
\item El plan puede ser sencillo: operar sólo dos o tres tipos de entradas en la gráfica de 5\~min, hacer \textit{scalping} con la mitad de la posición y \textit{swing} con la otra mitad usando un \textit{break‑even stop}.
\item La disciplina para \textbf{seguir el plan} es crítica; unas pocas fallas pueden anular semanas de beneficio.
\item Brooks advierte: ``\textit{No conviertas el mercado en un casino}''; el mercado castiga la falta de planificaci\'on.
\end{itemize}

%----------------------------------------------------------------------%
\begin{center}
  \begin{minipage}{0.9\textwidth}
    \justifying
    {\LARGE \textbf{PLAN DE TRADING (v.\,3)}}\\[2ex]   % ← solo \\[2ex]
    {\itshape
    «Cada entrada refleja una intenci\'on;\\
    cada salida, una consecuencia.\\
    Reconocer el \textsc{volumen}, respetar la \textsc{clavicular}\\
    y esperar el \textsc{secado} tras el \textsc{testeo}\\
    me conducir\'a a decisiones objetivas.\\
    S\'olo la }\textsc{disciplina}{\itshape\ convierte las ideas en resultados.»}
  \end{minipage}
\end{center}


%--------------------------------------------------------------------%
\section\*{Objetivos}

\begin{enumerate}
  \item \textbf{1\textsuperscript{o} Objetivo}.\\
        Alcanzar un \textsc{ratio aciertos/fallos} $\ge 60/40$ con
        \textsc{riesgo/beneficio} medio $\le 1{:}2$ durante un periodo de
        \textbf{12 meses} en operativa real.
  \item \textbf{2\textsuperscript{o} Objetivo}.\\
        Beneficio neto diario medio de \textbf{25~\$}, operando m\'aximo
        \textbf{3 horas} por sesi\'on.
  \item \textbf{3\textsuperscript{o} Objetivo}.\\
        Consolidar consistencia antes de \textbf{diciembre~2025};
        aumentar lotaje s\'olo tras tres cierres mensuales positivos.
\end{enumerate}


%--------------------------------------------------------------------%
\section\*{Focalizaci'on}

\subsection\*{Estrategia}
Trader \textbf{scalper e intrad'ia}: b'usqueda de impulsos de
\textit{precio+volumen} en marcos de 2\~min, con extensi'on a \emph{swing} intrad'ia si el contexto lo justifica.

\subsection\*{Horario operativo}
\begin{itemize}
\item \textbf{mini S\&P~~500}: 08:30--13:30.
Sin operativa los primeros 30~~min tras apertura.
\item Suspensi'on durante noticias de alto impacto hasta que el \textsc{volumen} vuelva a la normalidad.
\end{itemize}

\subsection\*{Temporalidad (\textit{Time Frame})}
\begin{itemize}
  \item Ejecuci\'on: \textbf{2\,min}.
  \item Contexto: 1~D $\rightarrow$ 240~min $\rightarrow$ 45~min.
\end{itemize}


\subsection\*{Direcci'on operativa}
Operar \emph{siempre} a favor de la tendencia vigente en la temporalidad de ejecuci'on.

%--------------------------------------------------------------------%
\section\*{Gesti'on del riesgo}

\begin{itemize}
  \item Cuenta de referencia: \textbf{1\,000~\$}.
  \item Riesgo por \textit{trade}: \textbf{2\,\%} (máx.\ 3\,\% en set-ups de alta probabilidad).
  \item Pérdida diaria máxima: \textbf{6\,\%}.\\
        Dos días consecutivos $\Rightarrow$ pausa y retroanálisis.
\end{itemize}


%--------------------------------------------------------------------%
\section\*{Sistema de trading}

\begin{itemize}
\item Basado en \textbf{An'alisis T'ecnico de Precio y Volumen}.
El \textsc{fundamental} s'olo para calendario macro.
\item \textbf{Regla de oro}: «\emph{Precio \$+{}\$ Volumen}. Precio manda, volumen acompa\~na».
\end{itemize}

\subsection*{Set-ups (entradas)}

\begin{description}
  \item[Set-Up VCS:] (Volumen–Clavicular–Secado)%
    \begin{enumerate}
      \item \textsc{Intenci\'on} rompe la clavicular con volumen alto.
      \item \textsc{Testeo} regresa $\le 50$\,\% del impulso con volumen decreciente.
      \item \textsc{Secado}: dos velas consecutivas con volumen $\le 70$\,\% del promedio $\Rightarrow$ entrada en ruptura.
    \end{enumerate}
  \item[Set-Up Falla Clavicular:] Segundo testeo sin ruptura + secado extremo $\Rightarrow$ operar contra el intento.
  \item[Set-Up Tap\'on:] Vela de rango estrecho con \textbf{volumen pico} en zona de agotamiento $\Rightarrow$ giro r\'apido \textit{scalp}.
\end{description}


\subsection*{Protocolo operativo}

\begin{enumerate}
  \item Detectar zona relevante (\textsc{clavicular}, HVN, pico de volumen).
  \item Esperar \textsc{intenci\'on} que rompa/valide la zona.
  \item Exigir \textsc{testeo} con volumen decreciente.
  \item Confirmar \textsc{secado}; sin secado $\Rightarrow$ \textbf{no trade}.
  \item Gestionar parcialmente en $R \ge 1.5$; resto a \textit{break-even}.
\end{enumerate}


%--------------------------------------------------------------------%
\section\*{Control y seguimiento}

\begin{itemize}
\item \textbf{Bit'acora diaria}: captura + comentario + emoci'on.
\item \textbf{Revisi'on semanal}: \% de cumplimiento, ratio\~R.
\item \textbf{Auditor'ia mensual}: curva de capital y drawdown.
\end{itemize}

%--------------------------------------------------------------------%
\section\*{Gesti'on emocional}

\begin{enumerate}
  \item Rutina \emph{pre-market}: respiraci\'on y visualizaci\'on.
  \item \textbf{Stop emocional}: 2 errores disciplinarios $\Rightarrow$ sesi\'on terminada.
  \item \textbf{Recompensa}: semana sin errores $\Rightarrow$ 10\,\% del beneficio para ocio.
\end{enumerate}


%--------------------------------------------------------------------%
\section\*{Gesti'on del dinero}

\begin{itemize}
  \item Beneficios: 40\,\% reinversi\'on,\ 40\,\% fondo de liquidez,\ 20\,\% diversificaci\'on.
  \item Revisi\'on trimestral seg\'un volatilidad y \textit{drawdown}.
\end{itemize}


%--------------------------------------------------------------------%
\section\*{Cl'ausula de mejora continua}

\begin{itemize}
  \item Revisar fiabilidad de \textit{set-ups}\ ($\ge 30$ muestras).
  \item Verificar secado efectivo.
  \item Ajustar horarios y activos por estacionalidad.
\end{itemize}


\vspace{1ex}
\begin{center}
\textit{«Operar es el arte de esperar cuando el mercado muestra su verdad con el \textsc{volumen}.»}
\end{center}

\newpage
%----------------------------------------------------------------------%
\section\*{Plantilla de Bit\'acora Diaria}

% Datos generales
\begin{tabularx}{\textwidth}{@{}lX@{}}
\textbf{Fecha:} & \\
\textbf{Sesi\'on:} & Londres / Nueva York / Asia\\
\textbf{Balance previo (\$):} & \\
\textbf{Objetivo del d\'ia (\$):} & \\
\textbf{P\'erdida m\'ax. (\%):} & 6,\% (seg\'un plan)\\
\textbf{Estado emocional inicial:} & \\
\end{tabularx}

% Tabla operaciones
\smallskip
\renewcommand{\arraystretch}{1.1}
\setlength{\tabcolsep}{4pt}
\begin{tabularx}{\textwidth}{@{}>{\raggedright\arraybackslash}p{1.3cm}ccccc>{\centering\arraybackslash}p{0.8cm}p{0.8cm}p{0.8cm}X@{}}
\toprule
Instr. & Dir. & Set-up & Entrada & SL & TP & H.E. & H.S. & R & Comentario \\
\midrule
& & & & & & & & & \\
& & & & & & & & & \\
\midrule
\multicolumn{6}{@{}r}{\textbf{Resultado neto (R):}} & \multicolumn{4}{l}{\hspace{1em}}\\
\bottomrule
\end{tabularx}

% Checklist
\smallskip
\begin{tabularx}{\textwidth}{@{}lccX@{}}
\toprule
\textbf{\'Item} & \textbf{S\'i} & \textbf{No} & \textbf{Observaciones}\\
\midrule
¿Se sigui\'o la rutina \textit{pre‑market}? & \qquad & \qquad & \\
¿Entradas con Intenci\'on–Test–Secado? & & & \\
¿Se respetaron SL y TP? & & & \\
¿Sesi\'on cerrada tras 2 errores? & & & \\
\midrule
\textbf{Aciertos (\%)} & & & \\
\textbf{Expectancia (R)} & & & \\
\midrule
\multicolumn{4}{@{}>{\raggedright\arraybackslash}p{\linewidth}@{}}{\textbf{Notas:}\newline
\-- Principales lecciones del d\'ia.\newline
\-- Ajustes para la sesi\'on siguiente.}\\
\bottomrule
\end{tabularx}

\newpage
% Cómo usar la plantilla
\smallskip
\section\*{C\'omo usar la plantilla}
\begin{itemize}
\item \textbf{Rellena} los datos generales al inicio.
\item \textbf{Registra} cada trade: Dir. = largo/corto; Set-up = VCS/Falla/Tap\'on.
\item \textbf{Completa} el checklist y m\'etricas al cierre.
\item \textbf{Guarda} capturas en la carpeta indicada.
\item \textbf{Revisa} la bit\'acora durante la auditor\'ia semanal.
\end{itemize}
